\documentclass[12pt, a4paper]{article}
\usepackage[utf8]{inputenc}
\usepackage{hyperref}
\usepackage{geometry}
\geometry{margin=1in}
\usepackage{tocloft}
\usepackage{booktabs}
\usepackage{enumitem}

\title{\textbf{TA Allocator - System Architecture and Design Document}}
\author{Department of Aerospace Engineering \\ Indian Institute of Technology Madras}
\date{Academic Year 2026}

\begin{document}

\maketitle
\tableofcontents
\newpage

\section{Project Overview}
The \textbf{TA Allocator} is a web-based administrative application designed to automate the Teaching Assistant (TA) allocation and faculty coordination workflow for the Department of Aerospace Engineering at IIT Madras. 

Historically, TA allocations were handled via fragmented Google Sheets and manual email correspondence. The TA Allocator automates this operational lifecycle by providing:
\begin{enumerate}[noitemsep]
    \item Multi-format CSV data ingestion for Courses, Students, Faculty, and Form Preferences.
    \item Bidirectional integration with the \textbf{Google Forms API (v1)} to dynamically populate live dropdown questions with eligible courses and currently unassigned students.
    \item Automated email dispatch via \textbf{Gmail SMTP (JavaMailSender)} for professor preference requests and official TA allotment confirmations.
    \item A constraint-satisfaction allocation algorithm that respects faculty choices, course priorities, availability locks, and course quotas.
\end{enumerate}

\section{System Architecture}
The application adopts an enterprise \textbf{Model-View-Controller (MVC)} pattern implemented using Java 17 and the Spring Boot framework.

\begin{itemize}
    \item \textbf{Presentation Layer (View):} Server-side rendered HTML using \textbf{Thymeleaf}. The presentation layer provides a master Admin Dashboard (\texttt{dashboard.html}) for file uploads and global metric tracking, along with a per-bucket management interface (\texttt{group-manager.html}).
    \item \textbf{Web/Controller Layer:} Handled by \texttt{AdminController.java}, which maps HTTP GET and POST requests, coordinates file streaming via \texttt{MultipartFile}, and manages view redirection with user flash messages.
    \item \textbf{Business Service Layer:} Modular Spring services (\texttt{DataUploadService}, \texttt{AllocationService}, \texttt{EmailService}, and \texttt{GoogleFormsService}) encapsulating data processing, third-party API communication, and allocation algorithms.
    \item \textbf{Persistence Layer (Model/Data):} Built on \textbf{Spring Data JPA} and \textbf{Hibernate} backed by an in-memory \textbf{H2 Database} (\texttt{jdbc:h2:mem:tadb}). State is maintained during runtime with low query latency and zero external database installation overhead.
\end{itemize}

\subsection{External Integrations}
\begin{itemize}
    \item \textbf{Google Forms API (v1):} Communicates with Google Cloud endpoints via OAuth 2.0 Service Account credentials (\texttt{credentials.json}). Updates form items in a single atomic \texttt{BatchUpdateFormRequest}.
    \item \textbf{Google SMTP Relay:} Dispatches automated notifications via \texttt{smtp.gmail.com} on port 587 using TLS/STARTTLS encryption and dedicated 16-character Google App Passwords.
\end{itemize}

\section{Module and Class Specifications}

\subsection{Data Models (\texttt{com.iitm.aero.taallocator.model})}
\begin{itemize}
    \item \textbf{\texttt{Course}}: Encapsulates course ID (\texttt{courseNo}), course title, assigned faculty name, faculty email, priority score, allocation group ID (\texttt{groupId}), and allotted assistants (\texttt{ta1}, \texttt{ta2}).
    \item \textbf{\texttt{Student}}: Captures roll number (\texttt{rollNo}), student name, CGPA, degree program (B.Tech, M.Tech, Ph.D.), current assigned course, and runtime availability boolean (\texttt{isAvailable}).
    \item \textbf{\texttt{Faculty}}: Stores instructor directory records, mapping names to institutional email addresses and internal department codes.
    \item \textbf{\texttt{PreferenceResponse}}: Represents ingested Google Form responses, storing the submitted course, instructor name, and ordered student choices (\texttt{taPreference1} through \texttt{taPreference4}).
\end{itemize}

\subsection{Repository Layer (\texttt{com.iitm.aero.taallocator.repository})}
Leverages Spring Data JPA interfaces that automatically generate runtime SQL queries:
\begin{itemize}
    \item \texttt{CourseRepository}: Provides \texttt{findByGroupId(double groupId)} for scoped allocation rounds.
    \item \texttt{StudentRepository}: Provides \texttt{findByIsAvailableTrue()} to query only unallocated candidates.
    \item \texttt{FacultyRepository}: Provides \texttt{findByNameIgnoreCase(String name)} for resilient faculty directory lookups.
    \item \texttt{PreferenceRepository}: Stores and queries raw form submission records.
\end{itemize}

\subsection{Service Layer (\texttt{com.iitm.aero.taallocator.service})}

\subsubsection{\texttt{DataUploadService}}
Parses uploaded CSV files utilizing \texttt{OpenCSV (v5.9)}. Implements dynamic header normalization to insulate the application against varying column naming conventions across different department spreadsheets:
\begin{itemize}[noitemsep]
    \item Resolves course columns matching \texttt{id}, \texttt{courseno}, \texttt{name}, \texttt{faculty\_name}, \texttt{priority}, \texttt{group\_id}.
    \item Resolves student columns matching \texttt{id}, \texttt{roll}, \texttt{name}, \texttt{cgpa}, \texttt{program}, \texttt{available}.
    \item Sets initial student availability flags based on prior course assignments.
\end{itemize}

\subsubsection{\texttt{AllocationService}}
Executes the allocation algorithm for a target \texttt{groupId}:
\begin{enumerate}[noitemsep]
    \item Retrieves all courses associated with the target group.
    \item Fetches faculty preference responses matching the course code.
    \item Evaluates special conditions, such as faculty opting out with \textit{"I do not want TAs"}.
    \item Iterates sequentially through faculty choices (\texttt{Preference 1} to \texttt{Preference 4}).
    \item Extracts student roll numbers using normalized parsing (\texttt{RollNo: Name}).
    \item Validates student existence and active availability (\texttt{student.isAvailable() == true}).
    \item Assigns the student as \texttt{ta1} or \texttt{ta2}, sets \texttt{isAvailable = false}, and persists state updates to the database.
    \item Enforces a strict maximum ceiling of 2 TAs per course.
\end{enumerate}

\subsubsection{\texttt{GoogleFormsService}}
Automates synchronization with Google Forms:
\begin{enumerate}[noitemsep]
    \item Authenticates using \texttt{GoogleCredentials} loaded from \texttt{credentials.json} with the \texttt{FORMS\_BODY} scope.
    \item Retrieves the target form schema via \texttt{forms.get(formId)}.
    \item Formats all unassigned students into descriptive dropdown items (\texttt{RollNo: Name [CGPA] (Program)}), prepending standard options (\textit{"No preference"}, \textit{"I do not want TAs"}).
    \item Identifies target questions by matching titles (\texttt{Select Course} and \texttt{TA Preference}).
    \item Executes a \texttt{BatchUpdateFormRequest} to rewrite dropdown options directly on Google's cloud servers.
\end{enumerate}

\subsubsection{\texttt{EmailService}}
Manages transactional SMTP email delivery via Spring's \texttt{JavaMailSender}:
\begin{itemize}[noitemsep]
    \item \textbf{Preference Collection Phase}: Formats tailored emails to professors containing course titles, submission deadlines, and direct Google Form links.
    \item \textbf{Confirmation Phase}: Formats final allotment notifications detailing assigned TAs and prompting faculty-student kickoff meetings.
    \item \textbf{Email Resolution Fallback}: Resolves instructor emails via three-tier lookup: exact directory match, normalized whitespace match, or embedded course record email.
\end{itemize}

\section{Operational Workflow and Group Segmentation}
Courses are segmented by \textbf{Group ID} (e.g., Group 0 for core courses, Group 1 for laboratory sessions). This segmentation provides two primary benefits:
\begin{enumerate}
    \item \textbf{Prioritized Allocation}: Higher-priority classes can be allocated and locked first.
    \item \textbf{Dynamic Candidate Pooling}: Students remaining unallocated after Group 0 instantly form the available candidate pool for Group 1.
\end{enumerate}

\section{Verification and Testing Strategy}
\begin{itemize}
    \item \textbf{Unit and Component Testing}: Validates parsing resilience against malformed CSV inputs, boundary conditions (courses requesting no TAs), and roll number string sanitization.
    \item \textbf{Integration Testing}: Confirms transactional rollbacks, state synchronization in H2, and accurate dropdown construction in \texttt{GoogleFormsService}.
    \item \textbf{End-to-End Simulation}: Validated using pre-packaged test data sets (\texttt{courses\_test.csv}, \texttt{students\_test.csv}, \texttt{faculty\_test.csv}, \texttt{preferences\_test.csv}) verifying the complete upload-to-email lifecycle.
\end{itemize}

\section{Milestone Status}
\begin{table}[h]
\centering
\renewcommand{\arraystretch}{1.4}
\begin{tabular}{|p{0.18\textwidth}|p{0.25\textwidth}|p{0.15\textwidth}|p{0.32\textwidth}|}
\hline
\textbf{Phase} & \textbf{Milestone} & \textbf{Status} & \textbf{Implementation Details} \\
\hline
\textbf{Phase 1} & Scaffolding \& UI & Completed & Spring Boot, Thymeleaf Dashboard, H2 DB setup. \\
\hline
\textbf{Phase 2} & CSV Ingestion & Completed & OpenCSV with resilient dynamic header mapping. \\
\hline
\textbf{Phase 3} & APIs \& SMTP & Completed & Google Forms API v1 client \& Gmail SMTP integration. \\
\hline
\textbf{Phase 4} & Allocation Core & Completed & Group-based matching, availability locking, max-2 limit. \\
\hline
\textbf{Phase 5} & Documentation & Completed & Detailed setup guide, README, and system design document. \\
\hline
\textbf{Phase 6} & Integration Testing & Completed & End-to-end verified with bundled test CSV datasets. \\
\hline
\textbf{Phase 7} & Production Rollout & Ready & Production deployment via executable JAR or container. \\
\hline
\end{tabular}
\end{table}

\end{document}
